% Curs 9, pagina 5 din 8 — transcriere a manuscrisului din 1994.
% Fragment destinat adăugării la sfârşitul lui courses/course-09.tex.
% Originalul nu a fost modificat. A se verifica notele editoriale.

\section*{Curs 9 - Pagina 5}

\noindent Not.\ cu $H=\langle a\rangle$, $\operatorname{ord}H=n$,
\[
H=\{e,a,\ldots,a^{n-1}\}.
\]

\noindent $b\notin H$.

\noindent Altfel $b=a^i$. Folosim faptul că $ba=a^{n-1}b$:
\[
a^{i+1}=a^{n+i-1}\ \big|\ \cdot a^{-i-1}
\quad\Longrightarrow\quad a^{n-2}=e\qquad\times
\]

\bigskip
\noindent $H\cap Hb=\varnothing\ \Longrightarrow\ \underbrace{H\cup Hb}_{2n\ \text{elem.}}\subseteq G$.

\bigskip
\noindent Dem.\ că $H\cup Hb=G$, adică
\[
G=\{e,a,a^2,\ldots,a^{n-1},\ b,ab,a^2b,\ldots,a^{n-1}b\}.
\]

\noindent Arăt că $H\cup Hb$ e parte stabilă. Cum e finită, va rez.\ că e subgr.

\[
a^i b\cdot a^s b=a^i a^{s(n-1)}b^2\in\langle a\rangle=H.
\]

\bigskip
\noindent Dacă $x\in G\ \Longrightarrow\ (\exists!)\,i,j$ cu
\[
0\leq i<n,\qquad 0\leq j<2,
\]
\noindent a.î.\ $x=a^i b^j$. Fie $y=a^s b^t$.\footnote{Exponenţii din a doua ramură a calculului următor sunt suprascrişi şi greu de descifrat în manuscris; $b^0$ este lectura provizorie, păstrată fără corecţie matematică. Este necesară verificarea cu originalul.}

\[
\begin{aligned}
xy&=a^i b^j a^s b^t\\
  &=\begin{cases}
      a^{i+s}b^t, & j=0,\\[2pt]
      a^{i+s(n-1)}b^0, & j=1.
    \end{cases}
\end{aligned}
\]
\bigskip
\noindent $f:G\longrightarrow\Delta_n$,\quad
$f(x)=R^iS^j\ \Longleftrightarrow\ x=a^i b^j$,
\[
0\leq i<n,\qquad 0\leq j<2.
\]
\noindent $f$ biject.

\noindent Arăt că $f(xy)=f(x)f(y)$.

\bigskip
\noindent\underline{\textbf{Acţiunea unui grup pe o mulţ.\ ($G$-mulţ.)}}

\bigskip
\noindent $(G,A,\psi)$,\quad $\psi:G\times A\longrightarrow A$ --- lege de comp.\ ext.\
pe $A$ cu operatori în $G$.

\noindent $G$ --- grup;\quad $A$ --- mulţ.\ nevidă.

\[
\left\{\begin{aligned}
(1)\quad &\psi(a,\psi(b,u))=\psi(ab,u), && a,b\in G,\ u\in A,\\
(2)\quad &\psi(e,u)=u.
\end{aligned}\right.
\]

\noindent Echivalent, scriind $\psi(a,u)=au$:
\[
\left\{\begin{aligned}
(1)\quad &a(bu)=(ab)u,\\
(2)\quad &eu=u.
\end{aligned}\right.
\]

\bigskip
\noindent\underline{Def.}\ Fie $G$ grup şi $A$ o mulţ.\ O aplic.\
$\psi:G\times A\longrightarrow A$ s.\ n.\ \underline{acţiune a gr.\ $G$ pe mulţ.\ $A$}
dacă avem satisfăcute (1) şi (2).
